% Options for packages loaded elsewhere
% Options for packages loaded elsewhere
\PassOptionsToPackage{unicode}{hyperref}
\PassOptionsToPackage{hyphens}{url}
\PassOptionsToPackage{dvipsnames,svgnames,x11names}{xcolor}
%
\documentclass[
  catalan,
  letterpaper,
  DIV=11,
  numbers=noendperiod]{scrartcl}
\usepackage{xcolor}
\usepackage[top=25mm,bottom=25mm,left=20mm,right=20mm]{geometry}
\usepackage{amsmath,amssymb}
\setcounter{secnumdepth}{5}
\usepackage{iftex}
\ifPDFTeX
  \usepackage[T1]{fontenc}
  \usepackage[utf8]{inputenc}
  \usepackage{textcomp} % provide euro and other symbols
\else % if luatex or xetex
  \usepackage{unicode-math} % this also loads fontspec
  \defaultfontfeatures{Scale=MatchLowercase}
  \defaultfontfeatures[\rmfamily]{Ligatures=TeX,Scale=1}
\fi
\usepackage{lmodern}
\ifPDFTeX\else
  % xetex/luatex font selection
\fi
% Use upquote if available, for straight quotes in verbatim environments
\IfFileExists{upquote.sty}{\usepackage{upquote}}{}
\IfFileExists{microtype.sty}{% use microtype if available
  \usepackage[]{microtype}
  \UseMicrotypeSet[protrusion]{basicmath} % disable protrusion for tt fonts
}{}
\makeatletter
\@ifundefined{KOMAClassName}{% if non-KOMA class
  \IfFileExists{parskip.sty}{%
    \usepackage{parskip}
  }{% else
    \setlength{\parindent}{0pt}
    \setlength{\parskip}{6pt plus 2pt minus 1pt}}
}{% if KOMA class
  \KOMAoptions{parskip=half}}
\makeatother
% Make \paragraph and \subparagraph free-standing
\makeatletter
\ifx\paragraph\undefined\else
  \let\oldparagraph\paragraph
  \renewcommand{\paragraph}{
    \@ifstar
      \xxxParagraphStar
      \xxxParagraphNoStar
  }
  \newcommand{\xxxParagraphStar}[1]{\oldparagraph*{#1}\mbox{}}
  \newcommand{\xxxParagraphNoStar}[1]{\oldparagraph{#1}\mbox{}}
\fi
\ifx\subparagraph\undefined\else
  \let\oldsubparagraph\subparagraph
  \renewcommand{\subparagraph}{
    \@ifstar
      \xxxSubParagraphStar
      \xxxSubParagraphNoStar
  }
  \newcommand{\xxxSubParagraphStar}[1]{\oldsubparagraph*{#1}\mbox{}}
  \newcommand{\xxxSubParagraphNoStar}[1]{\oldsubparagraph{#1}\mbox{}}
\fi
\makeatother

\usepackage{color}
\usepackage{fancyvrb}
\newcommand{\VerbBar}{|}
\newcommand{\VERB}{\Verb[commandchars=\\\{\}]}
\DefineVerbatimEnvironment{Highlighting}{Verbatim}{commandchars=\\\{\}}
% Add ',fontsize=\small' for more characters per line
\usepackage{framed}
\definecolor{shadecolor}{RGB}{241,243,245}
\newenvironment{Shaded}{\begin{snugshade}}{\end{snugshade}}
\newcommand{\AlertTok}[1]{\textcolor[rgb]{0.68,0.00,0.00}{#1}}
\newcommand{\AnnotationTok}[1]{\textcolor[rgb]{0.37,0.37,0.37}{#1}}
\newcommand{\AttributeTok}[1]{\textcolor[rgb]{0.40,0.45,0.13}{#1}}
\newcommand{\BaseNTok}[1]{\textcolor[rgb]{0.68,0.00,0.00}{#1}}
\newcommand{\BuiltInTok}[1]{\textcolor[rgb]{0.00,0.23,0.31}{#1}}
\newcommand{\CharTok}[1]{\textcolor[rgb]{0.13,0.47,0.30}{#1}}
\newcommand{\CommentTok}[1]{\textcolor[rgb]{0.37,0.37,0.37}{#1}}
\newcommand{\CommentVarTok}[1]{\textcolor[rgb]{0.37,0.37,0.37}{\textit{#1}}}
\newcommand{\ConstantTok}[1]{\textcolor[rgb]{0.56,0.35,0.01}{#1}}
\newcommand{\ControlFlowTok}[1]{\textcolor[rgb]{0.00,0.23,0.31}{\textbf{#1}}}
\newcommand{\DataTypeTok}[1]{\textcolor[rgb]{0.68,0.00,0.00}{#1}}
\newcommand{\DecValTok}[1]{\textcolor[rgb]{0.68,0.00,0.00}{#1}}
\newcommand{\DocumentationTok}[1]{\textcolor[rgb]{0.37,0.37,0.37}{\textit{#1}}}
\newcommand{\ErrorTok}[1]{\textcolor[rgb]{0.68,0.00,0.00}{#1}}
\newcommand{\ExtensionTok}[1]{\textcolor[rgb]{0.00,0.23,0.31}{#1}}
\newcommand{\FloatTok}[1]{\textcolor[rgb]{0.68,0.00,0.00}{#1}}
\newcommand{\FunctionTok}[1]{\textcolor[rgb]{0.28,0.35,0.67}{#1}}
\newcommand{\ImportTok}[1]{\textcolor[rgb]{0.00,0.46,0.62}{#1}}
\newcommand{\InformationTok}[1]{\textcolor[rgb]{0.37,0.37,0.37}{#1}}
\newcommand{\KeywordTok}[1]{\textcolor[rgb]{0.00,0.23,0.31}{\textbf{#1}}}
\newcommand{\NormalTok}[1]{\textcolor[rgb]{0.00,0.23,0.31}{#1}}
\newcommand{\OperatorTok}[1]{\textcolor[rgb]{0.37,0.37,0.37}{#1}}
\newcommand{\OtherTok}[1]{\textcolor[rgb]{0.00,0.23,0.31}{#1}}
\newcommand{\PreprocessorTok}[1]{\textcolor[rgb]{0.68,0.00,0.00}{#1}}
\newcommand{\RegionMarkerTok}[1]{\textcolor[rgb]{0.00,0.23,0.31}{#1}}
\newcommand{\SpecialCharTok}[1]{\textcolor[rgb]{0.37,0.37,0.37}{#1}}
\newcommand{\SpecialStringTok}[1]{\textcolor[rgb]{0.13,0.47,0.30}{#1}}
\newcommand{\StringTok}[1]{\textcolor[rgb]{0.13,0.47,0.30}{#1}}
\newcommand{\VariableTok}[1]{\textcolor[rgb]{0.07,0.07,0.07}{#1}}
\newcommand{\VerbatimStringTok}[1]{\textcolor[rgb]{0.13,0.47,0.30}{#1}}
\newcommand{\WarningTok}[1]{\textcolor[rgb]{0.37,0.37,0.37}{\textit{#1}}}

\usepackage{longtable,booktabs,array}
\newcounter{none} % for unnumbered tables
\usepackage{calc} % for calculating minipage widths
% Correct order of tables after \paragraph or \subparagraph
\usepackage{etoolbox}
\makeatletter
\patchcmd\longtable{\par}{\if@noskipsec\mbox{}\fi\par}{}{}
\makeatother
% Allow footnotes in longtable head/foot
\IfFileExists{footnotehyper.sty}{\usepackage{footnotehyper}}{\usepackage{footnote}}
\makesavenoteenv{longtable}
\usepackage{graphicx}
\makeatletter
\newsavebox\pandoc@box
\newcommand*\pandocbounded[1]{% scales image to fit in text height/width
  \sbox\pandoc@box{#1}%
  \Gscale@div\@tempa{\textheight}{\dimexpr\ht\pandoc@box+\dp\pandoc@box\relax}%
  \Gscale@div\@tempb{\linewidth}{\wd\pandoc@box}%
  \ifdim\@tempb\p@<\@tempa\p@\let\@tempa\@tempb\fi% select the smaller of both
  \ifdim\@tempa\p@<\p@\scalebox{\@tempa}{\usebox\pandoc@box}%
  \else\usebox{\pandoc@box}%
  \fi%
}
% Set default figure placement to htbp
\def\fps@figure{htbp}
\makeatother



\ifLuaTeX
\usepackage[bidi=basic,shorthands=off]{babel}
\else
\usepackage[bidi=default,shorthands=off]{babel}
\fi
\ifLuaTeX
  \usepackage{selnolig} % disable illegal ligatures
\fi


\setlength{\emergencystretch}{3em} % prevent overfull lines

\providecommand{\tightlist}{%
  \setlength{\itemsep}{0pt}\setlength{\parskip}{0pt}}



 


\KOMAoption{captions}{tableheading}
\makeatletter
\@ifpackageloaded{tcolorbox}{}{\usepackage[skins,breakable]{tcolorbox}}
\@ifpackageloaded{fontawesome5}{}{\usepackage{fontawesome5}}
\definecolor{quarto-callout-color}{HTML}{909090}
\definecolor{quarto-callout-note-color}{HTML}{0758E5}
\definecolor{quarto-callout-important-color}{HTML}{CC1914}
\definecolor{quarto-callout-warning-color}{HTML}{EB9113}
\definecolor{quarto-callout-tip-color}{HTML}{00A047}
\definecolor{quarto-callout-caution-color}{HTML}{FC5300}
\definecolor{quarto-callout-color-frame}{HTML}{acacac}
\definecolor{quarto-callout-note-color-frame}{HTML}{4582ec}
\definecolor{quarto-callout-important-color-frame}{HTML}{d9534f}
\definecolor{quarto-callout-warning-color-frame}{HTML}{f0ad4e}
\definecolor{quarto-callout-tip-color-frame}{HTML}{02b875}
\definecolor{quarto-callout-caution-color-frame}{HTML}{fd7e14}
\makeatother
\makeatletter
\@ifpackageloaded{caption}{}{\usepackage{caption}}
\AtBeginDocument{%
\ifdefined\contentsname
  \renewcommand*\contentsname{Taula de continguts}
\else
  \newcommand\contentsname{Taula de continguts}
\fi
\ifdefined\listfigurename
  \renewcommand*\listfigurename{Llista de figures}
\else
  \newcommand\listfigurename{Llista de figures}
\fi
\ifdefined\listtablename
  \renewcommand*\listtablename{Llista de taules}
\else
  \newcommand\listtablename{Llista de taules}
\fi
\ifdefined\figurename
  \renewcommand*\figurename{Figura}
\else
  \newcommand\figurename{Figura}
\fi
\ifdefined\tablename
  \renewcommand*\tablename{Taula}
\else
  \newcommand\tablename{Taula}
\fi
}
\@ifpackageloaded{float}{}{\usepackage{float}}
\floatstyle{ruled}
\@ifundefined{c@chapter}{\newfloat{codelisting}{h}{lop}}{\newfloat{codelisting}{h}{lop}[chapter]}
\floatname{codelisting}{Llistat}
\newcommand*\listoflistings{\listof{codelisting}{Llista de llistats}}
\makeatother
\makeatletter
\makeatother
\makeatletter
\@ifpackageloaded{caption}{}{\usepackage{caption}}
\@ifpackageloaded{subcaption}{}{\usepackage{subcaption}}
\makeatother
\usepackage{bookmark}
\IfFileExists{xurl.sty}{\usepackage{xurl}}{} % add URL line breaks if available
\urlstyle{same}
\hypersetup{
  pdftitle={Memòria de la Pràctica PPC: Implementació i Simulació d'un Controlador de Potència Activa en una Planta Fotovoltaica},
  pdfauthor={Arnau Coronado Nadal},
  pdflang={ca},
  colorlinks=true,
  linkcolor={blue},
  filecolor={Maroon},
  citecolor={Blue},
  urlcolor={Blue},
  pdfcreator={LaTeX via pandoc}}


\title{Memòria de la Pràctica PPC: Implementació i Simulació d'un
Controlador de Potència Activa en una Planta Fotovoltaica}
\usepackage{etoolbox}
\makeatletter
\providecommand{\subtitle}[1]{% add subtitle to \maketitle
  \apptocmd{\@title}{\par {\large #1 \par}}{}{}
}
\makeatother
\subtitle{Assignatura: GeDIT --- Control de Planta Fotovoltaica}
\author{Arnau Coronado Nadal}
\date{2026-11-12}
\begin{document}
\maketitle

\renewcommand*\contentsname{Taula de continguts}
{
\setcounter{tocdepth}{3}
\tableofcontents
}

\section{Introducció i Objectius}\label{introducciuxf3-i-objectius}

La ràpida transició energètica i la creixent penetració de la generació
renovable distribuïda en els sistemes elèctrics de potència exigeixen
transformar les plantes fotovoltaiques en unitats actives, gestionables
i capaces de proporcionar serveis de suport a la xarxa elèctrica. Les
plantes fotovoltaiques modernes ja no poden operar de manera passiva
injectant la màxima energia solar instantània, sinó que han de complir
estricts \textbf{codis de xarxa} (\emph{Grid Codes}) establerts pels
operadors dels sistemes de transport (TSO) i distribució (DSO).

Per coordinar l'operació simultània de desenes o centenars d'inversors
fotovoltaics distribuïts al parc i garantir el compliment de les ordres
de control al punt d'interconnexió (\emph{Point of Common Coupling},
PCC), es requereix un element de control centralitzat anomenat
\textbf{Power Plant Controller (PPC)}.

\begin{tcolorbox}[enhanced jigsaw, arc=.35mm, bottomrule=.15mm, bottomtitle=1mm, breakable, colback=white, colbacktitle=quarto-callout-note-color!10!white, colframe=quarto-callout-note-color-frame, coltitle=black, left=2mm, leftrule=.75mm, opacityback=0, opacitybacktitle=0.6, rightrule=.15mm, title=\textcolor{quarto-callout-note-color}{\faInfo}\hspace{0.5em}{Objectius Principals de la Pràctica}, titlerule=0mm, toprule=.15mm, toptitle=1mm]

\begin{itemize}
\tightlist
\item
  \textbf{Comprendre la funció del Power Plant Controller (PPC)} com a
  cervell de control centralitzat en plantes de generació renovable de
  gran escala.
\item
  \textbf{Modelar una planta fotovoltaica simplificada a
  MATLAB/Simulink}, representant la dinàmica de resposta dels inversors,
  les restriccions d'assombrament atmosfèric i els retards puros del
  sistema de comunicacions.
\item
  \textbf{Implementar l'algorisme de control de potència activa
  mitjançant una M-S-Function a MATLAB}, avaluant la sintonització de
  paràmetres del controlador PI (\(K_p\), \(K_i\)), la velocitat de
  resposta i la gestió de fenòmens de saturació de l'integrador
  (\emph{windup}).
\item
  \textbf{Avaluar el control combinat amb emmagatzematge d'energia
  (BESS)} per a la limitació estricta de rampa de potència i
  l'optimització de l'estat de càrrega (\emph{State of Charge},
  \(SOC\)).
\end{itemize}

\end{tcolorbox}

\begin{center}\rule{0.5\linewidth}{0.5pt}\end{center}

\section{Marc Normatiu i Funcionalitats del Power Plant Controller
(PPC)}\label{marc-normatiu-i-funcionalitats-del-power-plant-controller-ppc}

\subsection{Marc Normatiu i Codis de
Xarxa}\label{marc-normatiu-i-codis-de-xarxa}

Els codis de xarxa recullen les especificacions tècniques i els límits
operatius obligatoris per mantenir la integritat, estabilitat i qualitat
de subministrament de la xarxa elèctrica. A Europa, la norma
\emph{Requirements for Generators} (RfG) desenvolupada per
\textbf{ENTSO-E} (Reglament UE 2016/631) categoritza els mòduls de parc
de potència (\emph{Power Park Modules}, PPM) i imposa requisits
d'observabilitat i controlabilitat a partir de \(800    ext{ W}\) de
potència nominal.

El incompliment dels marges reglamentaris en les variables elèctriques
desencadena severos problemes operatius:

\begin{itemize}
\tightlist
\item
  \textbf{Desviacions de Freqüència (Variable Global):} Directament
  vinculades a l'equilibri entre la potència activa generada i la
  demanda. Sobrepassar els límits provoca variacions no desitjades en la
  velocitat de motors síncrons industrials, descalibratge de filtres
  passius d'harmònics, alteració en bancs de condensadors i saturació
  del nucli magnètic en transformadors i inductàncies, incrementant les
  pèrdues per corrents de Foucault i distorsionant l'ona.
\item
  \textbf{Desviacions de Tensió (Variable Local):} Relacionades amb els
  fluxos de potència reactiva a les xarxes de transport i distribució.
  Les sobretensions eleven el risc de perforació en els aïllaments dels
  equips i saturen màquines elèctriques, mentre que les subtensions
  incrementen la corrent i les pèrdues per efecte Joule, limitant la
  capacitat d'evacuació dels generadors i escurçant la vida útil de les
  càrregues electròniques.
\end{itemize}

\subsection{Funcions Principals del
PPC}\label{funcions-principals-del-ppc}

El PPC centralitza les lectures en temps real de l'equip de mesura al
PCC (\(P\), \(Q\), \(V\), \(f\)) i processa les consignes externes
rebudes del TSO/DSO mitjançant la infraestructura SCADA/SmartBridge.
Executa els seus algorismes en cicles periòdics amb temps de resposta
que oscil·len des de \(100   ext{ ms}\) fins a pocs segons.

\begin{tcolorbox}[enhanced jigsaw, arc=.35mm, bottomrule=.15mm, bottomtitle=1mm, breakable, colback=white, colbacktitle=quarto-callout-important-color!10!white, colframe=quarto-callout-important-color-frame, coltitle=black, left=2mm, leftrule=.75mm, opacityback=0, opacitybacktitle=0.6, rightrule=.15mm, title=\textcolor{quarto-callout-important-color}{\faExclamation}\hspace{0.5em}{Funcions Regulades del PPC}, titlerule=0mm, toprule=.15mm, toptitle=1mm]

\begin{enumerate}
\def\labelenumi{\arabic{enumi}.}
\tightlist
\item
  \textbf{Limitació de Potència Activa (\emph{Active Power
  Curtailment}):} Ajust del sostre màxim de potència injectable
  (\(P_{TSO}\)) per gestionar congestions a les línies d'evacuació o
  evitar sobrecàrregues a la xarxa.
\item
  \textbf{Control de Rampa (\emph{Ramp Rate Limitation}):} Limitació de
  les pendents màximes de variació de potència (\(\Delta P / \Delta t\))
  davant fluctuacions de radiació per pas de núvols.
\item
  \textbf{Regulació de Freqüència (\emph{f-P Droop Control}):} Modulació
  de la potència activa en resposta a desviacions de freqüència
  mitjançant característiques estàtiques de caiguda (\emph{droop} del
  2\% al 10\%).
\item
  \textbf{Control de Potència Reactiva, Factor de Potència i Tensió
  (\(V-Q\)):} Control independent de la corrent en quadratura (\(I_q\))
  per mantenir el factor de potència o la tensió al PCC.
\item
  \textbf{Suport davant Faltes (\emph{Fault Ride-Through - LVRT/HVRT}):}
  Injecció ultra-ràpida de corrent reactiva
  (\(\Delta I_q / I_N \ge 2.0    ext{ p.u.}\)) en menys de
  \(20    ext{ ms}\).
\end{enumerate}

\end{tcolorbox}

\begin{center}\rule{0.5\linewidth}{0.5pt}\end{center}

\section{Descripció del Model de la Planta
Fotovoltaica}\label{descripciuxf3-del-model-de-la-planta-fotovoltaica}

Per analitzar la dinàmica de control del PPC en ràgim de transitoris
``lents'', s'utilitza un model representatiu simplificat a
MATLAB/Simulink que prescindeix de les dinàmiques de commutació
electromagnètica ràpida (\(<20 ext{ ms}\)).

\subsection{Arquitectura de la Planta}\label{arquitectura-de-la-planta}

\begin{itemize}
\tightlist
\item
  \textbf{Potència nominal total de la planta (\(P_{plant}\)):}
  \(10    ext{ MW}\) (\(10000   ext{ kW}\)).
\item
  \textbf{Agrupació de generació:} 5 inversors fotovoltaics idèntics
  connectats en paral·lel, cadascun amb una potència nominal de
  \(P_{nom,i} = 2    ext{ MW}\) (\(2000    ext{ kW}\)).
\item
  \textbf{Punt d'Interconnexió (PCC):} Punt comú on se sumen totes les
  corrents/potències dels inversors per a la seva evacuació a la xarxa.
\end{itemize}

\subsection{Model Simplificat de l'Inversor
Fotovoltaic}\label{model-simplificat-de-linversor-fotovoltaic}

El conjunt format pel camp fotovoltaic i el seu corresponent inversor es
modela com un sistema de primer ordre amb un bloc de saturació a la
sortida que representa el límit físic per irradiància i temperatura
(\(P_{pv-av}\)):

\[P_{pv,i}(s) = rac{1}{    au_{pv}s + 1} \cdot P_{pv-set,i}^*(s)\]

On: * \(lpha \in [0, 1]\) és la consigna adimensional en per unitat
(p.u.) emesa pel PPC. * \(P_{pv-set,i}^* = lpha \cdot P_{nom,i}\) és la
consigna de potència activa en \$ ext\{kW\}\$ per a l'inversor \(i\). *
\$ au\_\{pv\} = 0.5 ext\{ s\}\$ (\(500  ext{ ms}\)) representa la
constant de temps equivalent de la resposta del bucle intern de
l'inversor. * \(P_{pv-av}\) és la potència màxima disponible al costat
DC sota les condicions solars instantànies.

\subsection{Retards de Comunicació}\label{retards-de-comunicaciuxf3}

La transmissió de dades entre la instrumentació del PCC, el PPC central
i les unitats de control local dels inversors introdueix retards de
xarxa. A Simulink, es representen mitjançant blocs de \textbf{retard
pur} (\emph{Comm delay}) fixats en \(20   ext{ ms}\)
(\(0.02    ext{ s}\)) tant al canal de realimentació de mesura com al
canal de consigna.

\begin{center}\rule{0.5\linewidth}{0.5pt}\end{center}

\section{Descripció de les Simulacions Realitzades i Model
PPC}\label{descripciuxf3-de-les-simulacions-realitzades-i-model-ppc}

\subsection{Algorisme del PPC en
M-S-Function}\label{algorisme-del-ppc-en-m-s-function}

El controlador central es programa mitjançant una \textbf{M-S-Function
de MATLAB} amb un període de mostratge discret de
\(T_{sample} = 200   ext{ ms}\) (\(0.2 ext{ s}\)).

L'algorisme llegeix en cada període l'error de potència activa al PCC:

\[e(k) = P_{pcc\_setpoint}(k) - P_{pcc\_meas}(k)\]

I calcula la nova consigna global \(lpha(k)\) mitjançant l'equació en
diferències d'un regulador PI:

\[P_{tot}(k) = P_{tot}(k-1) + K_p \cdot [e(k) - e(k-1)] + K_i \cdot T_{sample} \cdot e(k)\]

\[lpha(k) = rac{P_{tot}(k)}{P_{plant}}\]

Aquesta única senyal \(lpha(k)\) es transmet a tots els inversors perquè
cadascun calculi la seva referència local
\(P_{ref,i} = lpha \cdot P_{nom,i}\).

\begin{center}\rule{0.5\linewidth}{0.5pt}\end{center}

\section{Anàlisi i Interpretació de
Resultats}\label{anuxe0lisi-i-interpretaciuxf3-de-resultats}

\subsection{Exercici 1: Justificació de la Necessitat del Controlador
Central}\label{exercici-1-justificaciuxf3-de-la-necessitat-del-controlador-central}

\begin{longtable}[]{@{}
  >{\raggedright\arraybackslash}p{(\linewidth - 8\tabcolsep) * \real{0.2000}}
  >{\raggedright\arraybackslash}p{(\linewidth - 8\tabcolsep) * \real{0.2000}}
  >{\raggedright\arraybackslash}p{(\linewidth - 8\tabcolsep) * \real{0.2000}}
  >{\raggedright\arraybackslash}p{(\linewidth - 8\tabcolsep) * \real{0.2000}}
  >{\raggedright\arraybackslash}p{(\linewidth - 8\tabcolsep) * \real{0.2000}}@{}}
\caption{Resultats de la simulació en bucle obert (Exercici
1)}\label{tbl-ejercicio1}\tabularnewline
\toprule\noalign{}
\begin{minipage}[b]{\linewidth}\raggedright
Escenari d'Assombrament
\end{minipage} & \begin{minipage}[b]{\linewidth}\raggedright
Consigna TSO (\(P_{TSO}\))
\end{minipage} & \begin{minipage}[b]{\linewidth}\raggedright
\(P_{pv-av}\) Inversors 1 i 2
\end{minipage} & \begin{minipage}[b]{\linewidth}\raggedright
Potència Lliurada PCC
\end{minipage} & \begin{minipage}[b]{\linewidth}\raggedright
Error Permanent al PCC
\end{minipage} \\
\midrule\noalign{}
\endfirsthead
\toprule\noalign{}
\begin{minipage}[b]{\linewidth}\raggedright
Escenari d'Assombrament
\end{minipage} & \begin{minipage}[b]{\linewidth}\raggedright
Consigna TSO (\(P_{TSO}\))
\end{minipage} & \begin{minipage}[b]{\linewidth}\raggedright
\(P_{pv-av}\) Inversors 1 i 2
\end{minipage} & \begin{minipage}[b]{\linewidth}\raggedright
Potència Lliurada PCC
\end{minipage} & \begin{minipage}[b]{\linewidth}\raggedright
Error Permanent al PCC
\end{minipage} \\
\midrule\noalign{}
\endhead
\bottomrule\noalign{}
\endlastfoot
\textbf{Condició Nominal (Sense Núvols)} & \(10000  ext{ kW}\)
(\(lpha=1.0\)) & \(2000 ext{ kW}\) & \(10000  ext{ kW}\) &
\(0  ext{ kW}\) \\
\textbf{Assombrament Parcial} & \(10000 ext{ kW}\) (\(lpha=1.0\)) &
\(1500 ext{ kW}\) & \(9000   ext{ kW}\) & \(-1000  ext{ kW}\) \\
\textbf{Assombrament Sever} & \(10000   ext{ kW}\) (\(lpha=1.0\)) &
\(500  ext{ kW}\) & \(7000   ext{ kW}\) & \(-3000  ext{ kW}\) \\
\textbf{Assombrament Sever en Curtailment} & \(5000 ext{ kW}\)
(\(lpha=0.5\)) & \(500  ext{ kW}\) & \(4000   ext{ kW}\) &
\(-1000  ext{ kW}\) \\
\end{longtable}

\textbf{Anàlisi Tècnic:} Els inversors 1 i 2 saturen la seva producció
al seu límit tèrmic/solar instantani (\(1500   ext{ kW}\) o
\(500    ext{ kW}\)). Els inversors 3, 4 i 5, en rebre únicament el
valor \(lpha\) fixat per la consigna global, no incrementen la seva
generació per compensar el dèficit. \textbf{Es demostra que el control
en bucle obert no garanteix el compliment dels compromisos amb el TSO}.

\subsection{Exercici 2: Sintonització i Dinàmica del Controlador
PI}\label{exercici-2-sintonitzaciuxf3-i-dinuxe0mica-del-controlador-pi}

\begin{longtable}[]{@{}
  >{\raggedright\arraybackslash}p{(\linewidth - 10\tabcolsep) * \real{0.1667}}
  >{\raggedright\arraybackslash}p{(\linewidth - 10\tabcolsep) * \real{0.1667}}
  >{\raggedright\arraybackslash}p{(\linewidth - 10\tabcolsep) * \real{0.1667}}
  >{\raggedright\arraybackslash}p{(\linewidth - 10\tabcolsep) * \real{0.1667}}
  >{\raggedright\arraybackslash}p{(\linewidth - 10\tabcolsep) * \real{0.1667}}
  >{\raggedright\arraybackslash}p{(\linewidth - 10\tabcolsep) * \real{0.1667}}@{}}
\caption{Avalació de combinacions de sintonització PI (Exercici
2)}\label{tbl-ejercicio2}\tabularnewline
\toprule\noalign{}
\begin{minipage}[b]{\linewidth}\raggedright
Consigna \(P_{TSO}\) {[}kW{]}
\end{minipage} & \begin{minipage}[b]{\linewidth}\raggedright
\(P_{pv-av,1,2}\) {[}kW{]}
\end{minipage} & \begin{minipage}[b]{\linewidth}\raggedright
\(K_p\)
\end{minipage} & \begin{minipage}[b]{\linewidth}\raggedright
\(K_i\)
\end{minipage} & \begin{minipage}[b]{\linewidth}\raggedright
Assoleix Consigna?
\end{minipage} & \begin{minipage}[b]{\linewidth}\raggedright
Comportament Dinàmic i Sobrepic
\end{minipage} \\
\midrule\noalign{}
\endfirsthead
\toprule\noalign{}
\begin{minipage}[b]{\linewidth}\raggedright
Consigna \(P_{TSO}\) {[}kW{]}
\end{minipage} & \begin{minipage}[b]{\linewidth}\raggedright
\(P_{pv-av,1,2}\) {[}kW{]}
\end{minipage} & \begin{minipage}[b]{\linewidth}\raggedright
\(K_p\)
\end{minipage} & \begin{minipage}[b]{\linewidth}\raggedright
\(K_i\)
\end{minipage} & \begin{minipage}[b]{\linewidth}\raggedright
Assoleix Consigna?
\end{minipage} & \begin{minipage}[b]{\linewidth}\raggedright
Comportament Dinàmic i Sobrepic
\end{minipage} \\
\midrule\noalign{}
\endhead
\bottomrule\noalign{}
\endlastfoot
\(7000     o 5000\) & \(2000\) & \(0.5\) & \(0\) & \textbf{No}
(\(P_{pcc} pprox 5200   ext{ kW}\)) & Sense sobrepic, però amb error
permanent. \\
\(7000     o 5000\) & \(2000\) & \(5.0\) & \(0\) & \textbf{Gairebé}
(\(P_{pcc} pprox 5020  ext{ kW}\)) & Redueix l'error de ràgim permanent,
però oscil·la. \\
\(7000     o 5000\) & \(2000\) & \(0.5\) & \(0.1\) & \textbf{Sí} &
Resposta extremadament lenta (\(t_{asent} pprox 15  ext{ s}\)). \\
\(7000     o 5000\) & \(2000\) & \(0.5\) & \(1.0\) & \textbf{Sí} &
\textbf{Ajust Òptim:} Temps d'asentament ràpid (\(\sim 3  ext{ s}\)) i
nul sobrepic. \\
\(7000     o 5000\) & \(2000\) & \(0.5\) & \(10.0\) & \textbf{Sí} &
\textbf{Inestable:} Oscil·lacions severes per interacció amb retards. \\
\(7000     o 5000\) & \(500\) & \(0.5\) & \(1.0\) & \textbf{Sí} & Eleva
\(lpha\) perquè els Inv 3-5 compensin la manca dels Inv 1-2. \\
\end{longtable}

\subsection{Exercicis 3 i 4: Saturació i Lògica
Anti-Windup}\label{exercicis-3-i-4-saturaciuxf3-i-luxf2gica-anti-windup}

\begin{Shaded}
\begin{Highlighting}[]
\CommentTok{\% Lògica Anti{-}Windup (Clamping) a la S{-}Function del PPC}
\KeywordTok{if}\NormalTok{ (}\VariableTok{alpha\_unstat} \OperatorTok{\textgreater{}} \FloatTok{1.0} \OperatorTok{\&\&} \VariableTok{error} \OperatorTok{\textgreater{}} \FloatTok{0}\NormalTok{) }\OperatorTok{||}\NormalTok{ (}\VariableTok{alpha\_unstat} \OperatorTok{\textless{}} \FloatTok{0.0} \OperatorTok{\&\&} \VariableTok{error} \OperatorTok{\textless{}} \FloatTok{0}\NormalTok{)}
    \CommentTok{\% Es congela l\textquotesingle{}integrador si la sortida està saturada i l\textquotesingle{}error mantindria la saturació}
    \VariableTok{KI\_integrator} \OperatorTok{=} \VariableTok{KI\_integrator\_prev}\OperatorTok{;}
\KeywordTok{else}
    \VariableTok{KI\_integrator} \OperatorTok{=} \VariableTok{KI\_integrator\_prev} \OperatorTok{+} \VariableTok{Ki} \OperatorTok{*} \VariableTok{T\_sample} \OperatorTok{*} \VariableTok{error}\OperatorTok{;}
\KeywordTok{end}
\end{Highlighting}
\end{Shaded}

\subsection{Exercici 5: Control Combinat amb Emmagatzematge BESS i
Rampes}\label{exercici-5-control-combinat-amb-emmagatzematge-bess-i-rampes}

L'emmagatzematge en bateries aporta el suport de potència activa
necessari per donar compliment a les restriccions de rampa màxima
(\(10\% P_{nom}/  ext{min} = 1    ext{ MW}/   ext{min}\)): *
\textbf{Control Directe de Rampa vs Filtrat Pas Baix:} El control
directe activa la bateria exclusivament quan la derivada de potència
sobrepassa el límit (\(\Delta P / \Delta t > \Delta P_{max}\)), reduint
dràsticament els cicles de càrrega/descàrrega. * \textbf{Estratègies
d'Estat de Càrrega (\(SOC\)):} Operar amb un \(SOC^*\) dinàmic
proporcional a la generació actual manté la bateria preparada davant
possibles rampes de baixada o pujada.

\[  ext{SOC}^*(t) = 0.4 + rac{0.6-0.4}{P_{plant}} \cdot P_{pv-meas}(t)\]

\begin{center}\rule{0.5\linewidth}{0.5pt}\end{center}

\section{Conclusions}\label{conclusions}

\begin{enumerate}
\def\labelenumi{\arabic{enumi}.}
\tightlist
\item
  \textbf{Invalidesa del Bucle Obert:} El control sense realimentació al
  PCC no pot compensar pertorbacions solars locals ni assombraments
  parcials als inversors.
\item
  \textbf{Robustesa del Controlador PI:} L'ajust \(K_p = 0.5\) i
  \(K_i = 1.0\) garanteix una resposta ràpida (temps d'asentament
  \(\sim 3    ext{ s}\)) i estable davant dels retards de comunicació de
  \(20   ext{ ms}\).
\item
  \textbf{Necessitat de l'Anti-Windup:} La tècnica de \emph{clamping} és
  indispensable per evitar el congelament de la resposta del PPC quan la
  consigna del TSO supera la capacitat de generació real disponible.
\item
  \textbf{Sinergia amb BESS:} La coordinació del PPC amb bateries permet
  suavitzar les rampes de potència provocades per passos de núvols,
  protegint la xarxa elèctrica i optimitzant la vida útil de la bateria
  mitjançant un \(SOC^*\) dinàmic.
\end{enumerate}




\end{document}
